\chapter{Literature review}
\label{sec:literature-review}

\lipsum[1-4]\todo{Replace with actual intro text}

\section{Radar polarimetry}

Although the foundational ideas of radar polarimetry date back to the 1970s and can be considered a mature concept -- with the 1980s and 1990s representing a golden period of theoretical and experimental development for \gls{sar} and meteorology~\parencite{leePolarimetricRadarImaging2017} -- its application potential in the automotive industry began to emerge only in the previous decade.
In remote sensing, polarimetry has long been the standard for classifying terrain textures or hydrometeor shapes; however, transferring these techniques to the automotive domain presents unique challenges.
Unlike the far-field, high-altitude geometries of \gls{sar}, automotive radar operates at grazing incidence angles with significant multipath interaction, in the near-to-intermediate field, and against highly dynamic, non-cooperative targets.

Consequently, the field of polarimetric automotive radar remains in its infancy, lacking robust methods for classifying dynamic \glspl{vru} such as pedestrians and cyclists.
Nonetheless, it holds a compelling promise for achieving higher reliability and sophistication in sensors for \gls{adas} and autonomous driving.
This promise has captured the attention of major automotive companies and research institutes, prompting viability confirmations~\parencite{tillyRoadUserClassification2021} and proof-of-concept system implementations~\parencite{tintiFullyPolarimetricAutomotive2024}.
The rapidly increasing innovation in this area is further evidenced by the recent emergence of monographs focusing specifically on polarimetric radar for automotive applications~\parencite{visentinPolarimetricRadarAutomotive2019}.

Establishing a polarimetric framework for \gls{vru} classification, however, places stringent demands on the underlying radar architecture.
The theoretical requirements for high polarization purity and \gls{xpd} directly inform the selection of radiating element types and feeding integration platforms, as will be discussed in detail in~\cref{sec:frontend}.
Furthermore, the need for simultaneous acquisition of the full scattering matrix for dynamic scenes necessitates a sophisticated large-aperture \gls{mimo} configuration.
The design of this \gls{mimo} topology and its impact on virtual aperture and polarimetric diversity will be the focus of~\cref{sec:mimo-topologies}.
This introductory chapter therefore provides the theoretical and phenomenological motivation for the hardware and topological developments that follow.

The exposition given in this chapter follows standard references to electrodynamics and radar polarimetry, such as~\textcite{zangwillModernElectrodynamics2012} and~\textcite{leePolarimetricRadarImaging2017}.
Additionally, a specialized monograph on polarimetric radar for automotive applications by~\textcite{visentinPolarimetricRadarAutomotive2019}, is discussed to provide further context and detail.\todo{Add citations throughout the chapter.}


\subsection{Fundamentals of radar polarimetry}
\label{sec:polarimetry-fundamentals}

Electromagnetic waves can be decomposed into orthogonal linear, circular, or elliptical polarization states, each associated with a specific temporal evolution of the electric field vector.
In radar applications, polarization serves as an additional dimension for characterizing scattering mechanisms: targets may preserve, transform, or depolarize the incident wave depending on their geometry, surface material, roughness, and orientation.
These transformations provide valuable classification features that are absent in scalar radar measurements~\parencite{leePolarimetricRadarImaging2017}.

% The description of polarization typically relies on the Jones vector for coherent fields and the Stokes vector or coherency matrix for partially coherent and incoherent fields.
% At automotive millimetre-wave frequencies, the high coherence of \gls{fmcw} radars allows Jones and coherency representations to remain applicable.
% The polarization purity of the transmitted and received waves, however, is strongly influenced by antenna \gls{xpd}, \gls{pcb} anisotropies, and mutual coupling -- highlighting the need for the careful array design discussed in~\cref{sec:frontend}.

\paragraph{Time convention.}
As common in engineering monographs, this text treats the convention of time as \emph{positive}, meaning that a scalar wave $w(\xi,t)$ propagating in the $\xi$ direction is expressed as $\exp\[\ii(\omega t - k\xi)\]$.
This is in contrast to physics literature, which often adopts a \emph{negative time} convention, leading to $\exp\[\ii(k\xi - \omega t)\]$.


\subsubsection{Polarization of electromagnetic waves}

While the full propagation of electromagnetic energy is governed by Maxwell's equations, for radar polarimetry it suffices to consider the solution for a monochromatic plane wave propagating in a direction given by the wave vector $\vec k = k\vec e_k$ with angular frequency $\omega = 2\pi f$.
Such electric and magnetic vector fields, $\vec E$ and $\vec B$, trace out an ellipse in the plane perpendicular to the direction of propagation $\vec k$, and can be expressed as~\parencite{zangwillModernElectrodynamics2012}
\begin{align}
    \vec E(\vec r, t) &= \vec E_\perp\exp\[\ii(\omega t - \vec k\cdot \vec r)\]
&
    &\text{ and }
&
    \vec B(\vec r, t) &= \vec B_\perp\exp\[\ii(\omega t - \vec k\cdot \vec r)\],
\end{align}
where $\vec E_\perp = E_1 \vec e_1 + E_2 \vec e_2$ and $c\vec B_\perp = \vec e_k \times \vec E_\perp$ are generally complex amplitude vectors lying in the plane perpendicular to $\vec k$, given by the orthonormal basis vectors $\vec e_1$ and $\vec e_2$.
The physical fields of a monochromatic plane wave are the \emph{real parts} of these expressions.
Focusing on the electric field component, we have
\begin{align}
    \vec E(\vec r,t) = \begin{bmatrix} |E_1|\exp(\ii\delta_1) \\ |E_2|\exp(\ii\delta_2) \end{bmatrix} \exp\[\ii(\omega t - \vec k\cdot \vec r)\],
\end{align}
where $\delta_1$ and $\delta_2$ denote the phase offsets of the orthogonal components $E_1$ and $E_2$, respectively.
The choice of the transverse basis vectors is arbitrary; however, in radar applications, it is customary to select the horizontal and vertical directions with respect to the established coordinate system.

The polarization state -- the geometric locus traced by the tip of the electric field vector over time -- is fully described by the Jones vector $\vec E_J$.
For a fully coherent wave, the \emph{Jones vector} is defined by the amplitudes $|E_1|$ and $|E_2|$ and the relative phase difference $\delta \coloneq \delta_2 - \delta_1$, taking the form
\begin{align}
    \label{eq:jones-vector}
    \vec E_J \coloneq \vec E(\mathbf o,0) = \begin{bmatrix} |E_1| \exp(\ii\delta_1) \\ |E_2| \exp(\ii\delta_2) \end{bmatrix} = \exp(\ii\delta_1) \begin{bmatrix} |E_1| \\ |E_2| \exp(\ii\delta) \end{bmatrix}.
\end{align}
%\index{Jones vector}
This representation is critical for the \gls{mimo} system design discussed in~\cref{sec:mimo-topologies}, as the transmitter and receiver chains operate coherently.
However, the Jones vector is strictly valid only for fully polarized, monochromatic waves.


\subsubsection{The Sinclair scattering matrix}
\label{sec:sinclair-matrix}

When concerned with scattering off targets, the relationship between the incident and scattered electric fields is of primary interest.
As discussed in the previous sections, polarization of a monochromatic plane wave at a given instant can be fully characterized by the Jones vector.
Furthermore, a set of two orthogonal Jones vectors forms a polarization basis.%1
    \todo{Add a section on selection of polarization basis. Note the choices used by various authors.}
Therefore, it is possible to establish a linear scattering model that relates the incident and scattered Jones vectors, $\vec E^{\mathrm{in}}$ and $\vec E^{\mathrm{sc}}$, respectively, via a complex scattering matrix:
\begin{align}
    \label{eq:sinclair-matrix}
    \vec E^{\mathrm{sc}} = \frac{\e^{-\ii kr}}{r} \vec S \vec E^{\mathrm{in}} = \frac{\e^{-\ii kr}}{r} \begin{bmatrix} S_{11} & S_{12} \\ S_{21} & S_{22} \end{bmatrix} \vec E^{\mathrm{in}}.
\end{align}
%\index{Sinclair scattering matrix}
The matrix $\vec S$, called the \emph{Sinclair scattering matrix}, encapsulates the target's polarimetric response.
In this representation, the diagonal elements correspond to the \emph{co-polar} scattering amplitudes, where transmit and receive polarizations are identical, while the off-diagonal elements represent the \emph{cross-polar} terms, describing the conversion of energy between orthogonal polarization states.
Furthermore, under the assumption of reciprocity in a linear, time-invariant medium, the scattering matrix is symmetric, yielding
\begin{align}
    \vec S = \exp(\ii\phi_{11}) \begin{bmatrix}
        |S_{11}| & |S_{12}|\exp\[\ii(\phi_{12} - \phi_{11})\] \\
        |S_{12}|\exp\[\ii(\phi_{12} - \phi_{11})\] & |S_{22}|\exp\[\ii(\phi_{22} - \phi_{11})\]
    \end{bmatrix}.
\end{align}
The absolute phase term $\exp(\ii\phi_{11})$ is not considered an independent parameter since it represents an arbitrary value given by the target range.
The main consequence of this symmetry is a reduction in the number of independent parameters from eight to five.%
    \footnote{
        Even after this reduction, fully polarimetric systems, capable of measuring the full scattering response, dispose of five independent parameters per resolution cell.
        This is in contrast to single-polarized systems, which measure only two, and it shows the increased complexity and information content of polarimetric measurements.
    }
However, practical considerations in automotive radar often challenge this ideal model by introducing near-field effects: Reciprocity holds for plane waves; for targets in the near-field of the array, such as a pedestrian situated \qty{2}{m} away from the antenna, the wavefront curvature may introduce deviations.

\paragraph{Scattering coordinate frameworks.}
When defining the polarimetric scattering matrix, it is necessary to assume a frame in which the polarization is defined.
Generally, there are two principal conventions: the \emph{\gls{fsa}} and the \emph{\gls{bsa}}.
While the \gls{fsa} convention, sometimes called \emph{wave-oriented}, defines the polarization basis for both incident and scattered waves so that the Cartesian $z$-axis always faces the $\vec k$ direction, the \gls{bsa} system operates by defining the basis of the scattered wave with respect to the receiving antenna.
This text assumes the \gls{bsa} convention, as is standard for monostatic radar.
This choice simplifies the scattering definition by defining a fixed coordinate system for both the incident and backscattered waves relative to the antenna.

\paragraph{From theoretical scattering to observed signatures.}
While the Sinclair matrix $\vec S$ provides a deterministic description of target scattering in an ideal environment, the transition to practical automotive sensing introduces two significant layers of complexity: hardware-induced distortion and the stochastic nature of distributed targets.

In practical scenarios, the measured matrix $\vec M$ is a transformation of the true scattering matrix $\vec S$ through the system's transfer functions.
This is typically modelled as%
    \footnote{
        In polarimetric calibration,~\cref{eq:measurement-model} is sometimes expressed \enquote{backwards}; that is, by equating the ideal scattering matrix to the transformed measurement.
        This framing is common because calibration texts often define $\vec R$ and $\vec T$ as the \emph{correction} matrices.
        Conversely, \cref{eq:measurement-model} expresses the parasitic effects in the \emph{forward} direction to emphasize the measurement process, hence taking $\vec R$ and $\vec T$ as the \emph{distortion} matrices.
    }
\begin{equation}
    \label{eq:measurement-model}
    \vec M = \vec R \vec S \vec T + \vec C + \vec N,
\end{equation}
where $\vec T$ and $\vec R$ represent the polarimetric imbalances of the transmitter and receiver chains, $\vec C$ accounts for antenna mutual coupling and leakage, and $\vec N$ is the additive noise.
Characterizing these terms is the primary objective of the calibration routines detailed in~\cref{sec:calibration}.

Second, in the presence of complex, non-point-like targets such as \glspl{vru}, a single Sinclair matrix is often insufficient to capture the depolarization caused by multiple scattering centres.
This necessitates the use of the second-order statistics introduced in the following section.


\subsubsection{Stokes parameters and the Poincaré sphere}
\label{sec:stokes-poincare}

In dynamic automotive scenarios, electromagnetic waves typically interact with \emph{distributed scatterers} -- extended targets such as road surfaces, vegetation, vehicles, or tunnel walls that comprise numerous independent scattering centres within a single resolution cell.
Because these sub-reflectors contribute random phase and amplitude fluctuations, the resulting interference induces depolarization, rendering the reflected wave partially polarized or incoherent.
To characterize these complex fields, the Stokes parameters are employed; they provide a phase-agnostic description of the wave's polarization state based on observable, time-averaged power measurements.

The transition from Jones formalism to Stokes parameters involves considering the Jones vector components as random processes, $E_1(t)$ and $E_2(t)$.
Taking the time-averaged%
    \footnote{
        Taking the outer product of a single Jones vector $\vec E_J=[E_1,E_2]^\T$ without averaging would yield a rank-1 matrix, corresponding to the theoretical ideal of a fully polarized, perfectly coherent wave.
        Temporal averaging is essential to capture partial polarization effects.
    }
outer product of the Jones vector yields the Hermitian positive semidefinite wave covariance matrix $\vec J$, often called the \emph{coherency matrix}:
\begin{align}
    \label{eq:coherency-matrix}
    \vec J = \langle \vec E_J \cdot \vec E_J^\dagger \rangle = \begin{bmatrix}
        \langle |E_1|^2 \rangle & \langle E_1 E_2^* \rangle \\
        \langle E_2 E_1^* \rangle & \langle |E_2|^2 \rangle
    \end{bmatrix},
\end{align}
%\index{coherency matrix}
where $\langle \cdot \rangle$ denotes temporal averaging.
Analogous to~\cref{eq:sinclair-matrix}, the diagonal elements represent the intensities of the orthogonal polarization components, while the off-diagonal elements capture the complex cross-correlation between the cross-polarized channels.

To facilitate convenient description through matrix decomposition, group theory is often employed.
Specifically, the formalism considers the $\mathrm{SU}(2)$ group basis consisting of the Pauli matrices:
\begin{align}
    \label{eq:pauli-matrices}
    \vec\sigma_0 = \begin{bmatrix*}[r] 1 & 0 \\ 0 & 1 \end{bmatrix*}, \quad
    \vec\sigma_1 = \begin{bmatrix*}[r] 1 & 0 \\ 0 & -1 \end{bmatrix*}, \quad
    \vec\sigma_2 = \begin{bmatrix*}[r] 0 & 1 \\ 1 & 0 \end{bmatrix*}, \quad
    \vec\sigma_3 = \begin{bmatrix*}[r] 0 & -\ii \\ \ii & 0 \end{bmatrix*}.
\end{align}
These matrices form a basis for decomposing the coherency matrix -- a decomposition technique commonly known as the \emph{Pauli decomposition}:
\begin{align}
    \label{eq:pauli-decomposition}
    \vec J = \frac{1}{2} \sum_{i=0}^3 g_i \vec\sigma_i = \frac 12 \begin{bmatrix}
        g_0 + g_1 & g_2 - \ii g_3 \\
        g_2 + \ii g_3 & g_0 - g_1
    \end{bmatrix}.
\end{align}
%\index{Pauli decomposition}
Here, the coefficients $g_i$ are the \emph{Stokes parameters}, defined as $g_i = \tr(\vec J \vec\sigma_i)$.
Together, they form the \emph{Stokes vector} $\vec g$, expressed as
\begin{align}
    \label{eq:stokes-vector}
    \vec g = \begin{bmatrix} g_0 \\ g_1 \\ g_2 \\ g_3 \end{bmatrix} = \begin{bmatrix}
        \langle |E_1|^2 \rangle + \langle |E_2|^2 \rangle \\
        \langle |E_1|^2 \rangle - \langle |E_2|^2 \rangle \\
        2 \Re\langle E_1 E_2^* \rangle \\
        -2 \Im\langle E_1 E_2^* \rangle
    \end{bmatrix}.
\end{align}
%\index{Stokes parameters}

As evident from~\cref{eq:stokes-vector}, the Stokes parameters are \emph{real-valued power quantities} directly measurable via standard \gls{rf} detectors.
Geometrically, the parameters $g_1, g_2, g_3$ span a three-dimensional orthogonal basis, mapping the polarization state onto the \emph{Poincaré sphere}, as illustrated in~\cref{fig:poincare-sphere}.
Within this topological framework, the parameter $g_0$ represents the total wave intensity and corresponds to the radius of the sphere.

Consequently, the condition of physical realizability -- derived from the positive semi-definiteness of the coherency matrix -- requires that the state vector lies either on the surface or within the volume of the sphere:
\begin{equation}
    \label{eq:stokes-inequality}
    g_0^2 \geq g_1^2 + g_2^2 + g_3^2.
\end{equation}
%\index{Poincar\'e sphere}
The equality in~\cref{eq:stokes-inequality} holds strictly for fully polarized waves, which map to the sphere's surface.
Conversely, the strict inequality characterizes partially polarized waves, which occupy the interior volume and are typical of clutter and distributed targets.
This geometric distinction naturally leads to the definition of the \emph{\gls{dop}} as the normalized radial distance from the origin:
\begin{align}
    \text{DOP} \coloneq \frac{\sqrt{g_1^2 + g_2^2 + g_3^2}}{g_0} = \sqrt{1-4\frac{\det(\vec J)}{\tr(\vec J)^2}}.
\end{align}
%\index{degree of polarization (DOP)}
This metric serves as a key discriminant between stable \gls{vru} scatterers and distributed clutter, a feature that will be exploited in later chapters~\cref{sec:processing} for target classification.
By applying \emph{polarimetric decompositions} to the coherency matrix $\vec J$, the scattering process can be decoupled into constituent mechanisms, such as \emph{single-bounce}, \emph{double-bounce}, and \emph{volume scattering}.
In the following sections, these decomposition theorems are utilized to extract robust polarimetric signatures, providing the feature set required for the classification models developed later in this work.

\begin{figure}[t]
    \centering
    \begin{tikzpicture}[scale=2, >=Stealth]
        \def\R{1.5}
        \def\Re{0.38}
        \definecolor{linblue}{RGB}{40,120,200}
        \definecolor{circgreen}{RGB}{60,180,75}
        % Sphere outline and axes
        \draw[black, thick] (0,0) circle (\R);
        \draw[black, thick, <-, name path=g1] (-\R-0.3,-0.5) node[below=2pt] {$g_1$} -- (\R+0.3,0.5);
        \draw[black, thick, ->, name path=g2] (-\R-0.3,0.5) -- (\R+0.3,-0.5) node[above=2pt] {$g_2$};
        \draw[black, very thick, ->] (0,-\R-0.4) -- (0,\R+0.4) node[left=2pt] {$g_3$};
        
        % Linear polarization (blue equator)
        \draw[linblue, thick, name path=equator] (0,0) ellipse [x radius=\R, y radius=\Re];
        % Horizontal and vertial
        \path [name intersections={of=equator and g1, by={V, H}}];
        \coordinate (H) at (H);
        \begin{scope}[shift={(H)}]
            \fill[linblue] (0,0) circle (0.03);
            \draw[linblue, <->, thick] (-0.25,0) -- (0.25,0);
            \node[below=12pt, text width=2cm, align=center] at (0,0) {H};
        \end{scope}
        \coordinate (V) at (V);
        \begin{scope}[shift={(V)}]
            \fill[linblue] (0,0) circle (0.03);
            \draw[linblue, <->, thick] (0,-0.25) -- (0,0.25);
            \node[above=12pt, text width=2cm, align=center] at (0,0) {V};
        \end{scope}
        % Diagonal linear
        \path [name intersections={of=equator and g2, by={M, P}}];
        \coordinate (P) at (P);
        \begin{scope}[shift={(P)}]
            \fill[linblue] (0,0) circle (0.03);
            \draw[linblue, <->, thick] (-0.2,-0.2) -- (0.2,0.2);
            \node[below=12pt, text width=2cm, align=center] at (0,0) {$\ang{45}$};
        \end{scope}
        \coordinate (M) at (M);
        \begin{scope}[shift={(M)}]
            \fill[linblue] (0,0) circle (0.03);
            \draw[linblue, <->, thick] (0.2,-0.2) -- (-0.2,0.2);
            \node[above=12pt, text width=2cm, align=center] at (0,0) {$-\ang{45}$};
        \end{scope}

        % Circular polarization (green meridian)
        \draw[circgreen, thick] (0,0) ellipse [x radius=\Re, y radius=\R];
        \begin{scope}[shift={(0,\R)}]
            \fill[circgreen] (0,0) circle (0.03);
            \draw[circgreen, ->, thick] (-45:0.15) arc (-45:-325:0.15);
            \node[above right=6pt, text width=2cm, align=left] at (0,0) {RHC};
        \end{scope}
        \begin{scope}[shift={(0,-\R)}]
            \fill[circgreen] (0,0) circle (0.03);
            \draw[circgreen, ->, thick] (45:0.15) arc (45:325:0.15);
            \node[below right=6pt, text width=2cm, align=left] at (0,0) {LHC};
        \end{scope}
    \end{tikzpicture}
    \caption{
        Representation of the polarization state on the Poincar\'e sphere.
        Fully polarized waves lie on the surface ($\text{DOP} = 1$), while partially polarized states lie within the volume ($\text{DOP} < 1$).
        The axes $g_1, g_2, g_3$ correspond to linear, linear-diagonal, and circular polarizations, respectively.
    }
    \label{fig:poincare-sphere}
\end{figure}


\subsection{Polarimetric target decomposition}
\label{sec:target-decomposition}

The Sinclair and coherency matrices derived in~\cref{sec:polarimetry-fundamentals} contain the complete polarimetric information of a target.
However, in their raw matrix form, they offer little direct insight into the physical geometry of the scatterer.
\emph{Target decomposition} theorems aim to invert this relationship, expressing the measured matrix as a linear combination of canonical scattering mechanisms -- such as flat plates, dihedrals, or dipoles -- thereby extracting semantic features for classification.

These methods are broadly categorized into \emph{coherent decompositions}, which operate on the Sinclair scattering matrix of deterministic targets, and \emph{incoherent decompositions}, which operate on the second-order statistics of distributed targets.


\subsubsection{Coherent decomposition}
\label{subsec:coherent-decomposition}

For deterministic targets with negligible noise and spatial variation, such as a car chassis or a corner reflector, the scattering process is fully coherent.
As described by~\textcite{gaglioneKrogagerDecompositionPseudoZernike2014}, a general coherent polarimetric decomposition of the Sinclair matrix $\vec S$ can be expressed as a linear combination of $M$ elementary scattering mechanisms; that is,
\begin{align}
    \vec S = \sum_{m=1}^M c_m \vec S_m,
\end{align}
where $\vec S_m$ are the canonical scattering matrices, encoding the response of the $m$-th canonical object, and $c_m$ are generally complex coefficients, including both amplitude and phase information of each scattering mechanism.\todo{Add a section on canonical scatterers prior to this.}

The most established framework for interpreting radar targets is the \emph{Pauli decomposition}, which projects the Sinclair matrix onto a scaled basis of Pauli matrices, $\{\sqrt{2}\vec\sigma_i\}_{i=0}^3$, where $\vec\sigma_i$ are defined in~\cref{eq:pauli-matrices}.%
    \footnote{
        The scaling factor of $\sqrt{2}$ ensures that the Euclidean norm of the target vector matches the Frobenius norm of the scattering matrix, thereby satisfying the requirement for \enquote{total power invariance}.
        The same reasoning applies to the lexicographic basis discussed later.
    }
The resulting vector representation of the scattering matrix, $\vec S$, is referred to as the \emph{Pauli scattering vector}, defined as
\begin{align}
    \label{eq:pauli-vector-4d}
    \vec k_P \coloneq \frac{1}{\sqrt{2}} \begin{bmatrix}
        S_{11} + S_{22} \\
        S_{11} - S_{22} \\
        S_{12} + S_{21} \\
        \ii(S_{12} - S_{21})
    \end{bmatrix} = \begin{bmatrix} a \\ b \\ c \\ d \end{bmatrix}.
\end{align}
The primary advantage of this basis is that it provides a direct physical interpretation of elementary scattering mechanisms.
Consequently, the squared magnitude of each Pauli component quantifies the contribution of a specific canonical mechanism to the total \gls{rcs}.
Specifically:
\begin{itemize}
    \item \emph{Single-bounce:} $|a|^2/2$ corresponds to odd-bounce scattering, typically arising from spheres, flat plates, or trihedral reflectors.
    In an automotive context, this mechanism dominates returns from flat surfaces, such as walls or the rear of a vehicle.
    \item \emph{Double-bounce:} $|b|^2/2$ corresponds to even-bounce scattering derived from dihedral structures.
    This is commonly observed in the corner-like features of a car (e.g., window frames and side mirrors) or the ground-wall interaction of a curb.
    \item \emph{Cross-polar:} $|c|^2/2$ represents the cross-polarization energy induced by dihedrals or dipoles rotated by $\ang{45}$ around the radar line-of-sight.
    \item \emph{Asymmetric:} $|d|^2/2$ accounts for non-reciprocal scattering mechanisms with asymmetric depolarizing effects.
\end{itemize}

In strict monostatic configurations, the principle of reciprocity applies, theoretically forcing the fourth component to zero; practically, it will contain noise or system artefacts.
In the ideal scenario of $d=0$, the Pauli vector reduces to a three-dimensional representation:
\begin{align}
    \label{eq:pauli-vector-3d}
    \vec k_P = \frac{1}{\sqrt{2}} \begin{bmatrix}
        S_{11} + S_{22} \\
        S_{11} - S_{22} \\
        2S_{12}
    \end{bmatrix}.
\end{align}
However, in the \emph{quasi-monostatic} scenarios relevant to this work -- where transmit and receive antennas are closely spaced but not co-located -- the asymmetry term $d$ may contain non-negligible system noise or phase imbalances which must be accounted for during calibration.
\todo{\textcite{bouwmeesterClassificationDynamicVulnerable2025} point out phase sensitivity of Pauli's decomposition.}

\paragraph{Lexicographic basis.}
An alternative representation frequently encountered in the literature is the \emph{lexicographic basis}, which is a scaled canonical basis of $\C^{2 \times 2}$:
\begin{align}
    \label{eq:lexicographic-basis}
    \Phi_L = \left\{
        2\begin{bmatrix} 1 & 0 \\ 0 & 0 \end{bmatrix},
        2\begin{bmatrix} 0 & 1 \\ 0 & 0 \end{bmatrix},
        2\begin{bmatrix} 0 & 0 \\ 1 & 0 \end{bmatrix},
        2\begin{bmatrix} 0 & 0 \\ 0 & 1 \end{bmatrix}
    \right\}.
\end{align}
This basis orders the target vector elements as $\vec k_L = [S_{11}, S_{12}, S_{21}, S_{22}]^\T$.
While is the native format for many radar hardware interfaces and is convenient for system calibration, it lacks the direct physical interpretability of the Pauli basis.

\paragraph{Alternative decompositions.}
While Krogager and Cameron proposed alternative coherent decompositions -- such as Krogager's approach using circular polarization to decompose targets into sphere, diplane, and helix components%
    \footnote{
        This is analogous to the Pauli decomposition but employs a circular polarization basis, which can be advantageous for identifying certain target geometries.
        While the components share similar interpretations (single-bounce and double-bounce), the helix component is physically distinct from Pauli's \enquote{tilted-wall component}; it captures chiral or asymmetric scattering mechanisms that are not explicitly represented as a unique shape in the Pauli basis.
    }
-- the Pauli basis remains the standard for coherent preprocessing due to its orthogonality and computational efficiency.


\subsubsection{Incoherent decomposition}
\label{sec:incoherent-decomposition}

In the context of \gls{vru} classification, targets are rarely simple point scatterers.
A pedestrian, for instance, is a complex aggregate of limbs with varying orientations and materials, possibly moving within a single resolution cell, creating a \emph{distributed target}.
To analyse such targets, it is necessary to move from the coherent vector $\vec k_P$ to the second-order statistics represented by the \emph{Pauli coherency matrix} $\vec T$, defined as a statistically averaged outer product of the Pauli vector:%
    \footnote{
        The Pauli coherency matrix $\vec T$ should not be confused with the general coherency matrix $\vec J$ defined in~\cref{eq:coherency-matrix}, which is a $2 \times 2$ matrix representing the wave's polarization state.
        While both are coherency matrices, they serve different purposes: $\vec J$ characterizes the polarization of the electromagnetic wave itself, whereas $\vec T$ encapsulates the statistical scattering properties of the target as represented in the Pauli basis.
    }
\begin{align}
    \vec T = \langle \vec k_P \cdot \vec k_P^\dagger \rangle,
\end{align}
where $\langle\cdot\rangle$ denotes statistical averaging over an arbitrary dimension.%
    \footnote{
        In \gls{sar} polarimetry, this averaging is typically performed over multiple looks or spatial pixels to reduce speckle noise.
        In automotive radar, where real-time processing is essential, this averaging may be performed temporally over multiple pulses or spatially across adjacent range, Doppler, or angle-of-arrival cells.
    }
This matrix is Hermitian positive semidefinite and, as such, is always diagonalizable by a unitary matrix $\vec U$, formed by orthonormal eigenvectors of $\vec T$, and the diagonalized matrix $\vec D$ features non-negative real entries on its main diagonal, which are the eigenvalues of $\vec T$.

It is worth noting that $\vec T$ is, under unitary transformation, mathematically equivalent to any matrix formed by an outer product of the Sinclair scattering matrix, regardless of the decomposition basis used for vector representation.
This means that the lexicographic basis introduced in~\cref{eq:lexicographic-basis}, as well as the Krogager circular and Cameron bases discussed above, can all be employed, making the following analysis techniques broadly applicable.

The most established method for analysing the coherency matrix, originally developed by Cloude and Pottier in 1997 for \gls{sar} imaging and remote sensing, is the \emph{Cloude--Pottier decomposition}, which assumes that there is always a dominant average scattering mechanism in each cell.
To generate estimates of the average target scattering matrix parameters, this method employs the eigenvalue expansion; that is
\begin{align}
    \vec T = \sum_{i=1}^3 \lambda_i \vec e_i \vec e_i^\dagger,
\end{align}
where are the eigenvalues, sorted in descending order ($\lambda_1 \ge \lambda_2 \ge \lambda_3$), and $\vec e_i$ are the orthogonal eigenvectors.

\paragraph{Noise subspace reduction.}
In a practical measurement utilizing the full 4D Pauli vector (Eq.~\ref{eq:pauli-vector-4d}), the coherency matrix is $4 \times 4$ with four eigenvalues.
However, under the monostatic assumption, the physical signal subspace is rank-3.
As discussed by~\textcite{visentinPolarimetricRadarAutomotive2019}, the fourth eigenvalue $\lambda_4$ can be attributed to \gls{awgn} in the non-reciprocal channel.
To enhance estimation accuracy, a noise subtraction step is often applied:
\begin{align}
    \lambda_i' = \lambda_i - \lambda_4 \quad \text{for } i=1,2,3,
\end{align}
assuming $\lambda_4$ represents the noise floor $N$.
Following this correction, the analysis proceeds on the reduced $3 \times 3$ subspace.

The resulting unitary matrix $\vec U = \begin{bmatrix} \vec e_1 & \vec e_2 & \vec e_3 \end{bmatrix}$ contains eigenvectors parametrized by five angular degrees of freedom, specifically:
\begin{align}
    \label{eq:eigenvector-parametrization}
    \vec e_i = \begin{bmatrix}
        \cos(\alpha_i)\exp(\ii\epsilon_i) \\
        \sin(\alpha_i) \cos(\beta_i) \exp(\ii \delta_i) \\
        \sin(\alpha_i) \sin(\beta_i) \exp(\ii \gamma_i)
    \end{bmatrix}, \quad i=1,2,3.
\end{align}

Based on this eigen-decomposition, three key parameters are defined.
First, defined from the logarithmic probability of the eigenvalues, the \emph{polarimetric entropy} measures the \enquote{randomness} of the scattering process:
\begin{align}
    \label{eq:polarimetric-entropy}
    H = -\sum_{i=1}^3 P_i \log_3(P_i), \quad\text{ where }\quad P_i = \frac{\lambda'_i}{\sum_{k=1}^3 \lambda'_k}.
\end{align}
This definition delineates the scattering processes: $H=0$ indicates a fully deterministic, single dominant mechanism without loss of polarimetric information, such as an isotropic point target, while $H=1$ indicates random noise or fully developed volumetric scattering with complete depolarization effects, such as dense foliage or complex clutter.

Second, the angle $\alpha_i$ is defined using the angular parametrization of eigenvectors from~\cref{eq:eigenvector-parametrization} to identify the physical mechanism associated with the $i$-th eigenvalue.
This polarimetric classifier, known as the \emph{mean alpha angle}, is defined as
\begin{align}
    \label{eq:mean-alpha-angle}
    \bar{\alpha} = \sum_{i=1}^3 P_i \alpha_i.
\end{align}
This feature effectively categorizes scattering mechanisms into three primary types:
\begin{itemize}
    \item $\bar{\alpha} \approx \ang{0}$: isotropic surface (road, car body);
    \item $\bar{\alpha} \approx \ang{45}$: dipole/volume (vegetation, potentially limbs);
    \item $\bar{\alpha} \approx \ang{90}$: isotropic dihedral (ground-wheel, curb).
\end{itemize}
It should be noted that $\alpha$ strongly depends on the angle of incidence; flat surfaces tend to yield lower $\alpha$ values at normal incidence and higher values at grazing angles due to varying reflection coefficients.

Finally, the \emph{polarimetric anisotropy} in another parameter defined as an eigenvalue ratio, constructed as a complementary description to entropy.
It characterizes the relative importance of the second and third eigenvalues with respect to their descending arrangement:
\begin{align}
    A = \frac{\lambda_2 - \lambda_3}{\lambda_2 + \lambda_3}.
\end{align}
This metric ranges from $A=0$, indicating equal contributions from the second and third scattering mechanisms, to $A=1$, where the third mechanism is negligible relative to the second.
Polarimetric anisotropy thus describes the complexity of the non-dominant scattering, providing a means to differentiate between targets that may exhibit similar entropy values.

Anisotropy is particularly meaningful when the entropy is high ($H > 0.7$): in this regime, it distinguishes between targets with two significant scattering mechanisms (high $A$) and those characterized by fully random or isotropic scattering (low $A$).
When entropy is low, the second and third eigenvalues are typically too small for anisotropy to provide reliable information, limiting its interpretability in such cases.

\paragraph{The $H/\alpha$ plane.}
The joint analysis of polarimetric entropy $H$ and mean alpha angle $\bar{\alpha}$ provides a powerful two-dimensional feature space for target classification, commonly referred to as the $H/\alpha$ plane.
This representation allows for intuitive visualization and separation of different scattering mechanisms, especially in \gls{sar} polarimetry, based on their polarimetric characteristics.
    \todo{Add figure of H/alpha plane with typical target clusters.}


\subsubsection{Simple polarimetric descriptors}
\label{sec:simple-polarimetric-descriptors}

While the target decomposition theorems discussed in~\cref{subsec:coherent-decomposition} and~\cref{sec:incoherent-decomposition} offer rigorous physical interpretations of the scattering mechanisms, their computational complexity -- requiring eigen-decomposition of the $3\times3$ coherency matrix for every resolution cell -- can be prohibitive for real-time automotive applications constrained by embedded hardware resources.
Consequently, simple polarimetric descriptors -- originally developed for hydrometeor classification~\parencite{bringiPolarimetricDopplerWeather2001} based on power ratios and correlation coefficients -- can offer low-complexity proxies for target classification.
However, the translation of these metrics from meteorological bands (S/C/X-band) to automotive millimetre-wave frequencies (\qtyrange{77}{81}{GHz}) requires careful consideration of wavelength-scale roughness and system limitations.

\paragraph{Differential reflectivity.}
Perhaps the most fundamental polarimetric ratio, differential reflectivity quantifies the power imbalance between horizontal and vertical polarizations.
It is defined logarithmically as
\begin{align}
    \label{eq:differential-reflectivity}
    Z_\mathrm{dr} = 10\log_{10}\(\frac{\langle |S_{11}|^2 \rangle}{\langle |S_{22}|^2 \rangle}\).
\end{align}
In meteorology, this metric serves as a feature for classifying raindrop oblateness.
In the automotive context, $Z_\mathrm{dr}$ provides a measure of the target's geometric orientation.
Passenger vehicles, being predominantly horizontally oriented structures, typically exhibit positive $Z_\mathrm{dr}$.
Conversely, pedestrians lack this stable horizontal dominance; their scattering response at \qty{79}{GHz} is distributed and fluctuating due to the roughness of clothing and limb motion~\parencite{deepRadarCrosssectionsPedestrians2020}, often resulting in a mean $Z_\mathrm{dr}$ near $0\,\mathrm{dB}$.
Thus, $Z_\mathrm{dr}$ acts as a robust discriminator for separating vehicles from non-horizontal clutter.

\paragraph{Linear depolarization ratio.}
The \gls{ldr} measures the system's ability to detect cross-polarized energy relative to the co-polarized return.
It is defined as
\begin{align}
    \label{eq:linear-depolarization-ratio}
    \mathrm{LDR} = 10\log_{10}\(\frac{\langle |S_{21}|^2 \rangle}{\langle |S_{22}|^2 \rangle}\).
\end{align}
An ideal monostatic radar observing a sphere or a flat plate at normal incidence results in zero cross-polarization ($\mathrm{LDR} \to -\infty$).
In contrast, complex targets with intricate geometries, such as bicycles, induce significant depolarization due to multiple scattering and non-orthogonal structural components~\parencite{bouwmeesterClassificationDynamicVulnerable2025}.
This can result in elevated \gls{ldr} values compared to flat plates, though detecting this weak cross-polarized return requires high \gls{snr}.
Practically, however, the utility of this metric is bounded by the antenna system's cross-polarization isolation; if the antenna leakage exceeds the target's depolarization response, the \gls{ldr} signature becomes corrupted, limiting its use to high-performance sensor architectures.

\paragraph{Co-polar correlation coefficient.}
To assess the coherence between the horizontal and vertical scattering centres, the co-polar correlation coefficient is utilized:
\begin{align}
    \label{eq:copolar-correlation-coefficient}
    \rho_{12} = \frac{|\langle S_{11} S_{22}^* \rangle|}{\sqrt{\langle |S_{11}|^2 \rangle \langle |S_{22}|^2 \rangle}}.
\end{align}
This statistical descriptor ranges from 0 to 1 and is potentially the most robust metric for automotive scenes.
Man-made objects with stable phase centres, such as vehicles, poles, guardrails, typically exhibit $\rho_{12}$ close to 1.
In contrast, distributed volume scatterers -- such as bushes, grass, and tree canopies -- decorrelate the orthogonal channels significantly due to their random orientation and depth, yielding lower correlation values.
This could make $\rho_{12}$ an effective pre-filter for suppressing vegetation clutter, a common source of false positives in radar perception.

These descriptors can be computed efficiently for each resolution cell and used as input features to lightweight classification algorithms.
For example, objects exhibiting very large negative \gls{ldr} values alongside high $\rho_{12}$ are likely to correspond to specular reflections off the road surface and buildings, which can be filtered out early.

\begin{custombox}[orange]{To-do}
    Add the following metrics used by~\textcite{bouwmeesterClassificationDynamicVulnerable2025} for classification:
    \begin{itemize}
        \item Polarimetric power ratio
        \begin{align}
            Q_{xy} = \frac{\sum_{i=1}^N |S_{xy}^i|^2}{\sum_{i=1}^N |S_{xx}^i|^2 + |S_{xy}^i|^2 + |S_{yx}^i|^2 + |S_{yy}^i|^2},
        \end{align}
        where $N$ is the number of range-Doppler detections corresponding to a target within a single frame.

        \item Distribution of polarimetric ratios in the range-velocity spectrum:
        \begin{quote}
            \enquote{
                This can be done by analysing the polarimetric ratios of the individual detections.
                They are computed as the ratio of the squared magnitude of the considered scattering parameter divided by the squared magnitude of all four scattering parameters for just a single detection, in contrast to the target polarimetric ratio which is computed using all detections corresponding to a target within a frame.
            }
        \end{quote}
    \end{itemize}
\end{custombox}


\section{Front-end technologies}
\label{sec:frontend}

The antenna element serves as the fundamental physical interface of the radar sensor, defining the initial boundary conditions for signal fidelity.
In the context of polarimetric \gls{mimo} automotive radar operating in the \qtyrange{77}{81}{GHz} band, the antenna design is governed by a stringent conflict between electromagnetic purity and mass-producibility.
While the synthesis of large virtual apertures relies on the placement of these elements, the quality of the polarimetric information is determined by the individual element's ability to maintain orthogonal polarization states over a wide \gls{fov}.

Achieving high \gls{xpd} at millimetre-wave frequencies is notoriously difficult, particularly when constrained by the low-profile, cost-sensitive packaging requirements of the automotive industry.
The antenna must not only exhibit robust isolation and gain stability but also mitigate the effects of surface waves and mutual coupling, which can degrade the orthogonality of the \gls{mimo} channels.
Consequently, the choice of antenna technology is not merely a component selection but a system-level architectural decision that dictates the achievable dynamic range of the polarimetric radar.


\subsection{Design requirements and constraints}

The transition from standard automotive radar to fully polarimetric imaging imposes a specific set of performance metrics that narrow the field of viable antenna topologies.
The primary driver is polarization purity across the scan volume: whereas legacy systems prioritize gain, a polarimetric sensor requires an \gls{xpd} exceeding \qty{20}{dB} across a wide azimuth \gls{fov} (typically up to $\pm\ang{60}$)~\parencite{tintiFullyPolarimetricAutomotive2024}.
This is challenging because the geometric projection of polarization vectors' orthogonality naturally degrades at oblique angles~\parencite{zhaoOverviewPolarimetryApplication2025}.

Furthermore, these electromagnetic requirements must be reconciled with the realities of automotive integration.
The antenna must be compatible with standard multi-layer \gls{pcb} or package-level -- typically \gls{bga} or \gls{ewlb} -- manufacturing processes, usually precluding bulky metallic waveguide flanges or machined horns.
Furthermore, mutual coupling becomes a critical parameter in dense \gls{mimo} arrays; insufficient isolation between co-located orthogonal ports of dual-polarized antenna elements can lead to signal leakage that is mathematically indistinguishable from target depolarization.
Finally, thermal stability is non-negotiable; the phase centre and resonant frequency must remain stable under the harsh temperature cycling of an automotive environment to prevent calibration drift in the virtual array manifold~\parencite{harterErrorAnalysisSelfcalibration2015}.


\subsection{Overview of radar front-end technology}

The literature presents a diverse array of structural architectures proposed for \qty{79}{GHz} operation.
These can be broadly categorized into planar printed structures -- which dominate the current commercial landscape -- and substrate-integrated or waveguide-based solutions that offer performance enhancements at the cost of manufacturing complexity.


\subsubsection{Planar printed technology}

The \emph{microstrip patch antenna} remains the standard radiating element for commercial automotive radar due to its seamless integration with low-cost multi-layer \gls{pcb} processes~\parencite{waldschmidtAutomotiveRadarFirst2021}.
In polarimetric applications, dual-polarized operation is typically achieved using orthogonal feeds (aperture-coupled or probe-fed) or proximity-coupled patches~\parencite{vallecchiDesignDualpolarizedSeriesfed2005,acimovicDualPolarizedMicrostripPatch2008,dengSeriesFedDualPolarizedSingleLayer2019}.
While highly manufacturable, patch antennas suffer from intrinsic limitations at microwave frequencies: they are prone to high dielectric loss and surface wave excitation which degrade efficiency and coupling, and their narrow impedance bandwidth can be a bottleneck for wideband chirps~\parencite{zang77GHzFully2024}.
Moreover, maintaining high \gls{xpd} at wide angles is difficult due to the inevitable cross-polar radiation from the feed lines.
Microstrip stub antennas, which utilize shaped stubs for improved impedance bandwidth and polarization control, pose a compelling alternative to conventional series patch arrays for compact, low-profile implementations~\parencite{tinti45degLinearlyPolarized2023}.

Antenna elements based on \emph{slotted radiation}, including cavity-backed slots and slanted slot arrays, offer an alternative planar approach.
Slots naturally exhibit high polarization purity and can be more robust against mutual coupling than patches~\parencite{sunWidebandCavitySlottedWaveguide2024,liDesignImplementationSIW2019}.
However, they are intrinsically single-polarized, and hence, integrating dual-polarized slot arrays often requires multi-layer feed networks~\parencite{ferreiraCompactDualCircularlyPolarized2017} that increase board complexity, and their bidirectional radiation pattern usually necessitates a back-cavity or reflector, adding to the vertical profile.

Emerging in recent years, \emph{thin-film antennas} leverage advanced deposition techniques to realize ultra-thin radiating structures directly on the substrate~\parencite{khanHybridThinFilm2017}.
While thin-film designs can support dual-polarization and flexible integration, they are typically very lossy at millimetre-wave frequencies, which limits their efficiency and practical deployment in automotive radar systems.


\subsubsection{Integrated waveguide technology}

To overcome the loss mechanisms of microstrip lines, \emph{\gls{siw}} and \emph{\gls{ltcc}} technologies confine the electromagnetic field within a synthesized waveguide structure embedded in the dielectric.
\gls{siw} antennas exhibit significantly lower transmission losses and superior element isolation compared to microstrip~\parencite{saeidi-maneshLowCrossPolarizationHighIsolation2019}.
Their enclosed nature provides excellent shielding, making them robust against interference and temperature-induced deformations.
However, the requisite via-fencing consumes considerable board real estate, posing a challenge implementing dense \gls{mimo} lattices without sacrificing grating lobe performance.

\emph{\Gls{gw}} and \emph{\gls{rgw}} technologies represent a significant evolution in low-loss millimetre-wave design, addressing the assembly bottlenecks of traditional hollow waveguides.
By utilizing an \emph{\gls{ebg}} surface, such as the \enquote{bed of nails} reported by~\textcite{kildalDesignExperimentalVerification2011}, to suppress parallel-plate modes, \gls{gw} technology achieves the low loss of air-filled waveguides while eliminating the requirement for conductive electrical contact between waveguide layers.

This \enquote{contactless} characteristic creates a unique advantage for automotive radar: it relaxes the stringent mechanical flatness and assembly torque requirements that drive up the cost of traditional split-block waveguides.
While early iterations relied on costly \gls{cnc} milling, recent advancements have enabled a transition to metallized plastic injection moulding.
As detailed in industrial studies by~\textcite{huegel3DWaveguideMetallized2022,bencivenniHighVolumeManufacturingMetallized2023}, this approach allows for the rapid prototyping of complex layers via 3D printing before scaling to high-precision injection moulding for mass production.
This workflow effectively shifts the complexity from assembly to the mould-design phase, resulting in high-efficiency, air-filled waveguide arrays commercially viable for the harsh thermal and vibrational environments of automotive sensing~\parencite{yongOverviewRecentDevelopment2023}.


\subsubsection{Volumetric antennas}

While impractical for conformal automotive integration, \emph{horn antennas} and \glspl{dra} serve as important benchmarks, hence their use in polarimetric system demonstrators, such as those presented by~\textcite{tintiFullyPolarimetricAutomotive2024,trummerPerformanceAnalysis792017}.
Horns exhibit outstanding gain and polarization discrimination but are volumetrically incompatible with bumper-integrated sensors.
\Glspl{dra} offer wide bandwidth and high radiation efficiency by eliminating metallic losses~\parencite{kishkWidebandTruncatedTetrahedron2003}, yet they face challenges regarding mechanical robustness and the precision assembly required to mount 3D dielectric blocks onto a planar radar transceiver.

\paragraph{Launcher-in-package technology.}
In recent years, \gls{lip} technology has emerged as a promising solution to the traditional challenges of integrating waveguide antennas with \gls{mmic}.
By embedding the waveguide transition directly within the package substrate, \gls{lip} minimizes the number of \gls{rf} transitions, thereby reducing insertion loss and reflection typically associated with conventional chip-to-waveguide interfaces~\parencite{mohammedAdvancementsMmWaveTechnology2024}.
While still an emerging technology, this advancement enables more efficient coupling between the transceiver and the antenna, making waveguide-based solutions more viable for compact automotive radar systems.


\subsection{Comparative analysis of antenna technologies}

A critical review of the literature reveals that antenna performance is not determined by a single design choice, but rather by the interaction between two distinct layers: the \emph{radiating element} and the \emph{integration platform}.
In many reported designs, these two aspects are conflated.
To provide a clearer map of the design space, this analysis decouples these layers, evaluating them separately before synthesizing the state-of-the-art findings.


\subsubsection{Radiating element typologies}
The choice of radiating element dictates the intrinsic bandwidth, polarization discrimination, and radiation pattern of the sensor.
\Cref{tab:radiating_elements} categorizes the fundamental element types found in automotive radar literature.

While microstrip patches dominate commercial implementations due to their low profile, they are intrinsically narrowband and prone to surface wave excitation.
In contrast, volumetric radiators like horns and \glspl{dra} offer superior bandwidth and polarization stability but face severe integration penalties.
Slot radiators occupy a middle ground, offering high \gls{xpd} but requiring complex back-cavity structures to suppress bidirectional radiation.


\subsubsection{Integration platforms}
The integration platform defines the feeding structure of the antenna system: it determines insertion loss, isolation between \gls{mimo} channels, and thermal stability.
As shown in \cref{tab:integration_platforms}, the industry standard (microstrip/\gls{pcb}) trades performance for cost, while academic research heavily favours waveguide-based structures to maximize signal fidelity.

\begin{table}[ht]
    \centering
    \caption{Comparison of radiating element types for polarimetric radar}
    \label{tab:radiating_elements}
    \small
    \begin{tabular}{p{3.8cm} p{1.8cm} p{1.8cm} p{1.8cm} >{\footnotesize\raggedright\arraybackslash}p{3cm}}
        \toprule
        \textbf{Element type} & \textbf{BW (\%)} & \textbf{XPD (dB)} & \textbf{Profile} & {\small\textbf{References}} \\
        \midrule
        Edge-fed patch & $<5$ & 15--20 & Very low & \parencite{dashDesignSeriesFed2021,baur77GHzAutomotive2013,hambergerMixedCircularLinear2019} \\
        \addlinespace
        Aperture-coupled patch & 5--15 & 20--30 & Low & \parencite{acimovicDualPolarizedMicrostripPatch2008,puskely5GSIWBasedPhased2022,animHighGainMillimeterWavePatch2021,karamiDevelopingBroadbandMicrostrip2022a,saeidi-maneshLowCrossPolarizationHighIsolation2019} \\
        \addlinespace
        Microstrip stub & $<5$ & 20--25 & Very low & \parencite{tinti45degLinearlyPolarized2023,salucciSplineShapedMicrostripEdgeFed2024} \\
        \addlinespace
        Thin-film antenna & $\sim10$ & 20--30 & Very low & \parencite{khanHybridThinFilm2017} \\
        \addlinespace
        Slotted waveguide & $<5$ & 15--25 & Low & \parencite{hehenberger77GHzFMCWMIMO2020,bauer79GHzResonantLaminated2013} \\
        \addlinespace
        Cavity-backed slot & $<5$ & $>25$ & Moderate & \parencite{ferreiraCompactDualCircularlyPolarized2017,sunWidebandCavitySlottedWaveguide2024,mosalanejadMultiLayerPCBBowTie2020,liDesignImplementationSIW2019,yangMillimeterWaveDualPolarizedDifferentially2020} \\
        \addlinespace
        Horn waveguide & $>15$ & $>25$ & High (3D) & \parencite{trummerPerformanceAnalysis792017,tintiFullyPolarimetricAutomotive2024,zhaoSIWQuasiPyramidHorn2024} \\
        \addlinespace
        Dielectric resonator & 10--15 & $\sim20$ & High (3D) & \parencite{guoDualpolarizedDielectricResonator2003,kishkWidebandTruncatedTetrahedron2003} \\
        \bottomrule
    \end{tabular}
\end{table}

\begin{table}[ht]
    \centering
    \caption{Comparison of integration and feeding platforms}
    \label{tab:integration_platforms}
    \small
    \begin{tabular}{p{3.4cm} p{2cm} p{2cm} p{2cm} >{\footnotesize\raggedright\arraybackslash}p{3.2cm}}
        \toprule
        \textbf{Platform} & \textbf{Loss} & \textbf{Isolation} & \textbf{Complexity} & {\small\textbf{References}} \\
        \midrule
        Planar PCB & High & Low & Very low & \parencite{hambergerMixedCircularLinear2019,dashDesignSeriesFed2021,baur77GHzAutomotive2013,tinti45degLinearlyPolarized2023,salucciSplineShapedMicrostripEdgeFed2024} \\
        \addlinespace
        Multi-layer PCB & Moderate & Moderate & Moderate & \parencite{aliakbari79GHzMultilayer2022,mosalanejadMultiLayerPCBBowTie2020,saeidi-maneshLowCrossPolarizationHighIsolation2019} \\
        \addlinespace
        Metallic waveguide & Very low & Very high & Moderate & \parencite{tintiFullyPolarimetricAutomotive2024,trummerPerformanceAnalysis792017,schulwitzMillimeterWaveDualPolarized2006,sunWidebandCavitySlottedWaveguide2024} \\
        \addlinespace
        SIW & Moderate & High & Moderate & \parencite{puskely5GSIWBasedPhased2022,karamiDevelopingBroadbandMicrostrip2022a,hehenberger77GHzFMCWMIMO2020,zhaoSIWQuasiPyramidHorn2024,liDesignImplementationSIW2019,yangMillimeterWaveDualPolarizedDifferentially2020} \\
        \addlinespace
        LTCC & Low & High & High & \parencite{bauer79GHzResonantLaminated2013} \\
        \addlinespace
        Gap waveguide & Low & Very high & High & \parencite{renAutomotivePolarimetricRadar2024,zang77GHzFully2024} \\
        \bottomrule
    \end{tabular}
\end{table}


\subsection{Synthesis of state-of-the-art trends}

Analysing the intersection of these element and platform choices reveals a clear dichotomy in the current state of the art.

\paragraph{Industrial baseline.}
Despite their electromagnetic limitations, microstrip patch arrays remain the incumbent solution for mass-market automotive radar.
The manufacturing maturity of standard \gls{pcb} processes allows for the low-cost integration of complex series-fed or series-parallel networks directly alongside the \gls{mmic}~\parencite{hartnerReliabilityPerformanceWafer2019}.
However, for polarimetric applications, this convenience comes at a cost: standard patches struggle to maintain \gls{xpd} beyond $\pm\qty{30}{\degree}$ scan angles~\parencite{yangMillimeterWaveDualPolarizedDifferentially2020}, and the mutual coupling between closely spaced elements in a dense \gls{mimo} array introduces phase errors that degrade virtual aperture synthesis and the performance of signal processing algorithms, as explored by~\textcite{arnoldEffectAntennaMutual2019}.

\paragraph{Performance frontier.}
Recent academic literature identifies \gls{gw} and \gls{rgw} as the primary candidates for overcoming the inherent dielectric losses and surface-wave coupling of high-frequency \gls{pcb} technology.
By creating a prohibiting \gls{ebg}, \gls{gw} structures suppress parallel-plate modes entirely without requiring a galvanic connection between the waveguide layers.
This \enquote{contactless} property is a critical industrial enabler; it significantly relaxes the mechanical assembly tolerances that make traditional hollow waveguides cost-prohibitive for high-volume production.
Through the adoption of metallized injection moulding and high-precision plastic stamping, the manufacturing complexity is shifted from individual assembly to the mould-design phase, allowing for the low-cost mass production of air-filled structures.
Consequently, \gls{gw} technology, particularly in combination with the \gls{lip} technology, represents the emerging \enquote{gold standard} for next-generation \gls{mimo} arrays, providing the polarimetric fidelity and thermal robustness required for high-resolution automotive sensing within a commercially viable form factor~\parencite{renGapwaveguideAutomotiveImaging2023}.
However, a major limitation of \gls{gw} and \gls{rgw} technologies is their large physical footprint.
A single antenna element cell is typically much larger than its printed counterpart, especially due to the requisite mushroom structure fence needed to isolate parallel channels.
This substantial size penalty makes the integration of large-aperture \gls{mimo} arrays highly impractical, if not entirely infeasible, within the constrained dimensions of automotive sensors.

\paragraph{The compromise.}
\gls{siw} architectures effectively bridge the gap between planar and volumetric designs.
By synthesizing a dielectric-filled waveguide using periodic via fences, \gls{siw} provides the electromagnetic shielding and low-loss characteristics of metallic waveguides while retaining the manufacturability of standard substrate processes.
While 3D metal waveguides indeed offer unmatched electromagnetic fidelity and lower insertion loss, they are mechanically complex, heavy, and expensive to assemble.
\gls{siw} effectively acts as a bridge -- giving the shielded, low-radiation benefits of a closed waveguide but fully integrated into a compact, low-cost, planar \gls{pcb} footprint.

Regarding the viability of \gls{siw} at millimetre-wave frequencies, concerns are often raised regarding \gls{pcb} manufacturing tolerances, cost, and field quality.
However, beyond its superior electromagnetic properties (such as self-shielding and TM mode cancellation), this sensitivity is precisely why \gls{siw} is advantageous over other planar technologies like microstrip or stripline.
The key advantage lies in how manufacturing tolerances affect performance.
In microstrip or stripline structures, performance is highly sensitive to the chemical etching of extremely narrow traces, often only around \qty{150}{\um} wide.
A standard \qty{\pm25}{\um} etching tolerance represents an \qty{18}{\percent} variation, which drastically alters line impedance and causes severe reflections and return loss.
With \gls{siw}, the boundaries are instead defined by via fences.
This shifts the tolerance from chemical etching to mechanical drilling or laser positioning, which is inherently much more precise.
Furthermore, minor tolerances in via placement primarily affect the waveguide's cutoff frequency -- a parameter that is far less sensitive than microstrip or stripline impedance -- resulting in a much more robust, repeatable, and high-yield manufacturing process.
It is reassuring that \gls{siw} are already widely utilized in millimetre-wave applications; the manufacturing supply chain is mature, and standard advanced \gls{pcb} fabricators handle these tolerances on a regular basis.

The versatility of this platform is exemplified by~\textcite{zhaoSIWQuasiPyramidHorn2024}, who demonstrated that multi-layer \gls{pcb} stacks can be used to synthesize quasi-pyramidal horn antennas entirely within the substrate.

Furthermore, the layered nature of these architectures facilitates advanced hybrid topologies.
For instance, designs by~\textcite{puskely5GSIWBasedPhased2022,yangMillimeterWaveDualPolarizedDifferentially2020} combine the isolation benefits of an integrated waveguide feed with the beam-shaping flexibility of printed elements, using coupling slots in a sandwiched ground layer to excite parasitic patches on the surface.

While standard organic \gls{pcb} laminates provide a cost-effective baseline for these designs, the \gls{siw} concept can be extended into more specialized material systems such as \gls{ltcc}.
In such implementations, the \gls{siw} benefit of high-density integration is paired with the superior thermal stability and dimensional tolerance of ceramics.
This makes \gls{ltcc}-based \gls{siw} particularly attractive for highly integrated package-level antennas, albeit at a higher implementation cost compared to conventional organic stacks.


\subsection{Emerging paradigms}
\label{sec:emerging_paradigms}

Beyond the established categories of printed and waveguide-based elements, several exploratory technologies are expanding the design space for automotive radar.
These paradigms prioritize the synthesis of ideal radiation characteristics over traditional fabrication constraints.

\paragraph{Magneto-electric dipoles.}
Originally developed to overcome the narrow bandwidth of standard patch antennas, the magneto-electric dipole has recently been adapted for millimetre-wave radar to address the specific requirements of polarimetric fidelity.
By combining a planar electric dipole with a complementary magnetic dipole (typically synthesized via shorted patches or via-fenced slots), \glspl{med} function effectively as Huygens sources.

This complementary excitation yields two distinct advantages for automotive sensing.
First, it achieves wide impedance bandwidths (usually over \qty{20}{\percent}), easily covering the full \qtyrange{76}{81}{GHz} band required for high-resolution 4D imaging~\parencite{li60GHzDualPolarizedTwoDimensional2016}.
Second, and more critically for polarimetry, \glspl{med} exhibit nearly identical E- and H-plane radiation patterns with low cross-polarization.
This pattern symmetry ensures that the radar's \enquote{view} of a target is consistent regardless of polarization orientation, potentially simplifying the calibration matrices required for precise \gls{doa} estimation.

\paragraph{Metasurface-enhanced isolation.}
As \gls{mimo} arrays become denser to support higher angular resolution, mutual coupling via surface waves becomes a dominant source of error.
Designers are increasingly leveraging metasurfaces -- periodic sub-wavelength structures etched into the \gls{pcb} ground or top layer -- to manipulate these boundary conditions.

Rather than functioning as radiating elements themselves, these structures often act as \gls{ebg} filters or soft surfaces.
By presenting a high surface impedance to grazing waves, they effectively trap or dissipate surface currents between adjacent patch elements.
Recent implementations have demonstrated isolation improvements of \qty{10}{dB} to \qty{20}{dB} with negligible impact on the primary radiation pattern, offering a planar solution to the coupling problems that traditionally required metallic wall isolation~\parencite{mohamedPerfectIsolationPerformance2019}.

\paragraph{Additive manufacturing and integrated lenses.}
Recent advances in high-precision additive manufacturing are transforming the landscape of millimetre-wave antenna prototyping.
Unlike traditional machining or moulding, 3D printing enables the fabrication of complex waveguide structures with intricate internal geometries, such as non-standard transitions, that are otherwise impossible to manufacture.
While currently more suited to rapid prototyping than high-volume automotive deployment, this technology offers a unique pathway for validating complex volumetric designs before committing to expensive tooling.

A compelling application of this paradigm is the direct 3D printing of integrated lens antennas, where dielectric lenses, such as Luneburg or Maxwell fish-eye profiles, are fabricated as part of the antenna structure~\parencite{choSeriesFedIntegratedLens2024}.
This approach enables volumetric gain enhancement within a compact footprint, but remains limited by the loss tangent and surface roughness of available printable materials, which can introduce additional losses and phase errors at \qty{79}{GHz}.
While promising for prototyping and specialized applications, further advances in printable materials and post-processing are needed for widespread automotive deployment.


\section{MIMO array design}
\label{sec:mimo-topologies}

Existing \gls{mimo} array designs for automotive radar predominantly use linear or planar layouts optimized for angular resolution and cost-efficient implementation~\parencite{texasInstrumentsAwr2994EvaluationModule2024}.
\Gls{tdm} \gls{mimo} architectures dominate current commercial systems, typically leveraging 3--4 transmitters with 4--8 receivers to synthesize virtual arrays of 12--32 channels, possibly cascading multiple such modules~\parencite{texasInstrumentsAwrxCascadedRadar2020}.
Recent research explores sparse 2D \gls{mimo} arrays, co-prime and nested sparse layouts, and wide-aperture imaging arrays.
However, most published designs assume a single polarization, leaving the interaction between sparsity and polarization purity largely unexplored.
Only a few prototypes integrate dual-polarization at W-band, and these often suffer from phase-centre mismatches, reduced aperture efficiency and severe cross-polar coupling.
As a result, no established array topology exists that simultaneously optimizes spatial resolution, polarization isolation and multiplexing feasibility for a fully polarimetric automotive radar.


\subsection{Fundamentals of MIMO radar}

The concept of \gls{mimo} antenna systems is a well-established and extensively published area within telecommunications, where transmitting data over multiple uncorrelated signals results in additional information about the state of the communication channel, aggregating effects like multipath propagation and other types of signal fading or delays.
This information is leveraged to estimate the channel matrix which is then used to compensate for propagation effects to recover the transmitted data.

While the foundational idea of \gls{mimo} systems has been successfully carried over to radar technology, the problem formulation is adapted: in telecommunications jargon, the primary goal of radar is to precise estimate the channel matrix, which encapsulates the desired information about propagation effects, such as time delay, Doppler shift, and angle of arrival, allowing for extraction of target range, velocity, azimuth, and angular position~\parencite{blissMultipleinputMultipleoutputMIMO2003,forsytheMultipleinputMultipleoutputMIMO2004}.
The success of this communications-to-radar transformation has inspired extensive research, leading to key advancements and further classifications, such as the division of the general array concept into \emph{co-located} and \emph{distributed} \gls{mimo}, as extensively discussed in~\parencite{fishlerMIMORadarIdea2004,fishlerSpatialDiversityRadars2006}.

Lastly, it should be noted that the transfer of \gls{mimo} concepts from communications to radar must be undertaken with care, respecting and integrating the existing knowledge of traditional radar technology.
Uncritical adaptation can lead to flawed conclusions, as seen in some early distributed \gls{mimo} literature, where certain contributions were either redundant rediscoveries of multistatic radar concepts or contained incorrect statements, as notes by \textcite{chernyakConceptMIMORadar2010}.

From the standpoint of antenna array design, the key feature of \gls{mimo} radar is the creation of a \emph{virtual array} through the combination of multiple transmit and receive antennas.
Mathematically, if an array has $N_\mathrm{Tx}$ transmitters and $N_\mathrm{Rx}$ receivers, and each transmitter emits an orthogonal waveform, the received signals can be separated and associated with each transmit-receive pair.
This effectively synthesizes $N_\mathrm{Tx} \times N_\mathrm{Rx}$ virtual elements, whose positions are determined by the sum (or, in some cases, the difference) of the physical locations of the transmit and receive antennas~\parencite{holderSystematicMethodsSynthesis2023}.
The resulting virtual array, which can be obtained as a convolution of the physical transmit and receive element positions, can have a larger aperture and finer angular resolution than the physical arrays alone, but the spatial distribution of these virtual elements depends on both the physical layout and the multiplexing scheme.

\paragraph{Co-located and distributed MIMO radar.}
A crucial distinction exists between these two configurations.
co-located \gls{mimo} radar systems feature transmit and receive antennas placed in proximity, enabling coherent transmission and detection.
This coherence allows for the formation of a virtual array with enhanced angular resolution and supports advanced adaptive processing techniques.
In contrast, distributed \gls{mimo} radar systems employ widely separated antennas, a technique known from multi-static radar, emphasizing on spatial diversity gain.
This diversity is particularly beneficial for overcoming target scintillation (i.e. fluctuations in radar cross-section) and for improving detection performance in complex environments.

Notably, research was carried out to also explore hybrid approaches that combine the benefits of both co-located and distributed \gls{mimo}.
For example, \textcite{xuIterativeGeneralizedLikelihoodRatio2007} discuss systems that leverage both coherent processing and spatial diversity, highlighting the following advantages:
\begin{itemize}
    \item \emph{Spatial diversity:} widely-separated antennas in distributed \gls{mimo} configurations help mitigate target scintillation and improve detection reliability.
    \item \emph{Flexible transmit beampattern design:} co-located \gls{mimo} enables optimization of the transmit covariance matrix, allowing power to be focused in directions of interest while minimizing correlation of backscattered signals.
    This leads to significant improvements in adaptive processing techniques.
    \item \emph{Enhanced resolution and clutter rejection:} \gls{mimo} radar scheme can achieve higher spatial resolution and improved clutter suppression compared to conventional radar systems.
\end{itemize}

The concept of combining several widely separated subarrays with each subarray containing closely spaced antennas is nowadays often implemented via \emph{multi-aperture multiplexing} and has been successfully applied in \emph{cooperative automotive radars}, improving angular resolution and field of view~\parencite{liangCooperativeAutomotiveRadars2022}.

\paragraph{Array architectures.}
A central theme in \gls{mimo} radar design is the trade-off between spatial diversity, coherent processing gain, and hardware complexity, which is reflected in the three principal array architectures: fully diverse \gls{mimo}, phased arrays, and hybrid \gls{mimo} topologies.

In a fully diverse \gls{mimo} configuration, each antenna element operates independently, transmitting and receiving orthogonal waveforms.
This maximizes spatial diversity and enables robust target detection in complex environments.
Mathematically, the transmit covariance matrix in this case is full-rank, reflecting the linear independence of the transmitted signals.
However, this requires a separate \gls{rf} processing chain for each element, increasing hardware cost and complexity.

Phased arrays, in contrast, employ coherent waveforms across all elements, with a common signal distributed via a beamforming network.
This approach enables electronic beam steering and coherent processing gain, but sacrifices spatial diversity since all elements transmit the same signal.
The transmit covariance matrix is fully degenerate (rank 1), as all signals are linearly dependent.

Many modern systems adopt a hybrid approach, first introduced and called \enquote{phased-MIMO radar} by~\textcite{hassanienPhasedMIMORadarTradeoff2010}, which combines aspects of both fully diverse \gls{mimo} and phased arrays.
In these systems, groups of antenna elements form subarrays that transmit coherent waveforms, while different subarrays operate independently with orthogonal signals.
The resulting transmit covariance matrix possesses non-deficient rank, spanning the range from unity to full rank, depending on the degree of the so-called \emph{waveform diversity}.
This design balances the benefits of spatial diversity and coherent processing, while managing hardware complexity.
The optimization of subarray partitioning and waveform assignment typically leads to a difficult, non-convex optimization problem, as discussed in the research area of \emph{beampattern synthesis}.


\subsection{Topology synthesis}

The concept of \emph{MIMO topology} can be defined as a mapping from the physical arrangement of transmit and receive antennas, together with the chosen multiplexing scheme, such as time-, frequency-, and code-division, onto the resulting virtual array.
This mapping determines the positions and spacings of the virtual elements, which can be uniform or non-uniform depending on the physical geometry and the multiplexing strategy.
Non-uniform virtual element spacing can lead to grating lobes, ambiguities, or degraded performance, making the design of the physical and virtual topology a critical aspect of \gls{mimo} radar engineering.

\paragraph{Definitions \& figures of merit.}
The \gls{mimo} design problem can be formalized as follows: given a desired virtual array configuration (e.g. uniform linear array with specific aperture and element spacing), determine the physical transmit and receive antenna placements and the multiplexing scheme that will synthesize this virtual array.
Key figures of merit for evaluating \gls{mimo} topologies include:
\begin{itemize}
    \item \emph{Virtual aperture:} Set of unique transmitter-receiver channels in the resulting virtual array.
    The aperture is often evaluated as the overall size of the virtual array, corresponding directly to the resulting \emph{angular resolution}.
    \item \emph{\gls{fov}:} The angular region over which the array can reliably detect and resolve targets, determined by the physical and virtual array geometry and element spacing.
    \item \emph{\gls{dof}:} The number of independent channels available for signal processing, which influences the ability to perform tasks like \emph{direction-of-arrival estimation} and \emph{parameter identifiability}.
    \item \emph{Element spacing:} The spacing between virtual elements, affecting \emph{grating lobes} and \emph{ambiguity}, especially in sparse layout, where the Shannon-Nyquist spatial sampling criterion is deliberately violated.
    \item \emph{\gls{sll}:} The level of side-lobes in the virtual array's beampattern, impacting \emph{clutter rejection} and \emph{target detection}.
    \item \emph{Mutual coupling:} The interaction between physical elements, which can distort the intended virtual array response, leading to \emph{errors in angle estimation} or \emph{degraded detection performance}.
    \item \emph{Implementation complexity:} The practical feasibility of the physical layout and multiplexing scheme, considering hardware constraints.
\end{itemize}


\subsubsection{Application requirements}

The synthesis of \gls{mimo} topologies is fundamentally dictated by the sensing scenario, making application requirements the primary driver of array design.
For automotive radar, the ability to resolve targets in both azimuth and elevation is essential, necessitating two-dimensional (2D) aperture extension.
Modern front-looking automotive radar sensors, as described by~\textcite{waldschmidtAutomotiveRadarFirst2021}, require a wide azimuth \gls{fov} of approximately $\pm\ang{30}$ to $\pm\ang{60}$ and a moderate elevation \gls{fov} of about $\pm\ang{15}$ to $\pm\ang{30}$.
These requirements exceed what can be achieved with simple non-uniform elevation extensions of moderate-aperture linear arrays, necessitating employment of either a large element count or sparse array techniques.

However, synthesizing optimal 2D virtual arrays is challenging: the number of possible transmit-receive combinations increases rapidly, and the synthesis remains a non-deterministic optimization problem with no general closed-form solution for constructing a virtual array with ideal properties from a fixed set of physical elements.
This combinatorial complexity -- combined with practical constraints such as hardware cost, mutual coupling, and calibration -- makes 2D \gls{mimo} topology synthesis a challenging and active research area \parencite{zhugeStudyTwoDimensionalSparse2012}.

In automotive radar, practical constraints such as limited sensor footprint and integration requirements restrict \gls{mimo} implementations almost exclusively to co-located array configurations.
This spatial constraint is a systemic requirement imposed by the compact form factors and mounting locations typical of automotive platforms.
As a result, the theoretical advantages of distributed \gls{mimo} -- such as enhanced spatial diversity -- are generally unattainable in this context.

The signal processing benefits of co-located \gls{mimo} for automotive radar are now well-established.
These include the synthesis of large virtual arrays, improved angular resolution, and the ability to apply advanced adaptive processing techniques.
The literature provides a comprehensive treatment of these advantages, with concepts such as parameter identifiability \parencite{liParameterIdentifiabilityMIMO2007} and virtual aperture extension \parencite{liMIMORadarColocated2007} serving as key metrics for system evaluation.
The maturity of the field is further underscored by the availability of specialized reference texts dedicated to \gls{mimo} radar theory and practice \parencite{liMIMORadarSignal2025}.

Despite significant progress, two major challenges remain.
First, the analytical synthesis of optimal 2D virtual arrays is still underdeveloped, leaving most practical designs reliant on numerical optimization or empirical patterns.
Second, there is a pressing need to bridge the gap between electromagnetic phenomena—such as mutual coupling and parasitic effects—and signal processing algorithms, as these physical realities can introduce errors that degrade radar performance.
Addressing these challenges is essential for advancing both the theoretical and practical capabilities of automotive \gls{mimo} radar.


\subsubsection{Topology synthesis methods}

The synthesis of \gls{mimo} array topologies is a mature field offering a spectrum of design methodologies, spanning rigid analytical constructions and flexible numerical optimization techniques.
Analytical approaches, particularly those governing uniform linear arrays and uniform planar arrays, provide closed-form solutions and valuable insight into the mapping between physical and virtual array geometries, typically derived using the concept of the \emph{difference co-array}~\parencite{aminSparseArraysRadar2024}.
While these deterministic designs offer predictable phase centres and well-characterized point spread functions, they are often limited to specific configurations and may not generalize to constrained automotive scenarios.

For more intricate designs -- such as non-uniform or highly sparse arrays, or when practical constraints like packaging and mutual coupling must be considered -- numerical optimization becomes essential.
In these cases, the synthesis problem is formulated using an objective function that encodes key figures of merit, including virtual aperture, side-lobe level, and ambiguity~\parencite{ericAmbiguityCharacterizationArbitrary1998}.
Optimization algorithms employed in this context range from evolutionary heuristics (such as genetic algorithms, simulated annealing, and differential evolution) to modern convex relaxation methods, iteratively exploring the design space to reveal Pareto-optimal trade-offs between angular resolution and side-lobe suppression~\parencite{huanSASASuperResolutionAmbiguityFree2023}.
However, the resulting \enquote{exotic} geometries often lack translational invariance and introduce significant implementation challenges, particularly in terms of manifold calibration and mutual coupling compensation.
These practical difficulties can obscure the validation of novel signal processing chains, as most optimization-driven designs assume idealized, minimum-scattering antennas -- an assumption that does not hold in real-world applications and can lead to degraded processing algorithm performance~\parencite{arnoldEffectAntennaMutual2019}.

\paragraph{Topology classes for automotive radar.}
When considering \gls{mimo} topologies for automotive applications, it is crucial to recognize the distinction between theoretical sparse concepts and industrially deployable classes.

The most established class is the uniform linear array and its two-dimensional extension, the uniform planar array.
These arrays feature regularly spaced elements -- typically at half-wavelength intervals ($d=\lambda/2$) -- resulting in an ambiguity-free field of view and direct compatibility with standard \gls{fft} algorithms.
Their regularity is particularly advantageous for polarimetric research; it minimizes phase centre mismatch errors that can otherwise corrupt the delicate phase relationships between horizontal and vertical polarization channels.
While the scalability of uniform arrays is limited by the proportional increase in hardware cost, they provide the most controlled environment for isolating and characterizing Doppler-resolved polarimetric signatures.

To address the aperture limitations of uniform arrays, structured sparse arrays have emerged as the current state-of-the-art in deployed automotive imaging radar (e.g. cascaded chipsets).
Unlike the random or highly irregular sparse arrays found in theoretical literature, industrially supported sparse designs rely on specific, repeatable patterns to extend the virtual aperture.
These designs offer a pragmatic compromise: they achieve the improved angular resolution necessary for modern sensing while maintaining enough structural regularity to be reliably calibrated in mass production.

